% 自己完結した日本語原稿用プリアンブル。
% LuaLaTeX + luatexja を前提とし、外部スタイルファイルをコピーしない。
\usepackage[margin=25mm]{geometry}
\usepackage{amsmath}
\usepackage{amssymb}
\usepackage{bm}
\usepackage{graphicx}
\usepackage{booktabs}
\usepackage{multirow}
\usepackage{array}
\usepackage{tabularx}
\usepackage{enumitem}
\usepackage{placeins}
\usepackage{caption}
\usepackage{subcaption}
\usepackage{xcolor}
\usepackage{xspace}
\usepackage{tikz}
\usetikzlibrary{arrows.meta,positioning,calc,fit,backgrounds,shapes.geometric,patterns}

\usepackage[numbers,sort&compress]{natbib}
\usepackage[colorlinks=true,
            linkcolor=black,
            citecolor=blue!60!black,
            urlcolor=blue!60!black]{hyperref}
\usepackage[nameinlink,capitalise]{cleveref}

\captionsetup{labelfont=bf, font=small}
\captionsetup[table]{skip=4pt}
\setlength{\parindent}{1em}
\setlength{\parskip}{0.2em}
\setlength{\tabcolsep}{5pt}
\renewcommand{\arraystretch}{1.12}
\linespread{1.05}

% --- 図表で共有する配色 -------------------------------------------------
\definecolor{courtgreen}{HTML}{2E7D5B}
\definecolor{gsblue}{HTML}{2F6FA5}
\definecolor{poseorange}{HTML}{D9772B}
\definecolor{lossred}{HTML}{B94A48}
\definecolor{softgray}{HTML}{F0F2F4}
\definecolor{unmeasuredgray}{HTML}{6B6B6B}

% --- 図で共有する TikZ スタイル ----------------------------------------
\tikzset{
  block/.style={draw, rounded corners=2pt, align=center, fill=softgray,
    minimum height=7mm, inner sep=4pt, font=\sffamily\scriptsize},
  datablock/.style={block, draw=gsblue, fill=gsblue!10},
  geomblock/.style={block, draw=courtgreen, fill=courtgreen!10},
  modelblock/.style={block, draw=poseorange, fill=poseorange!10},
  objblock/.style={block, draw=lossred, fill=lossred!10},
  gateblock/.style={block, draw=black!45, fill=white,
    font=\sffamily\fontsize{6.4}{7.4}\selectfont},
  flow/.style={-{Latex[length=2.0mm,width=1.4mm]}, semithick},
  supflow/.style={flow, dashed, draw=lossred},
  boundary/.style={draw=gsblue, dashed, rounded corners=3pt, inner sep=3pt},
  figtitle/.style={font=\sffamily\bfseries\footnotesize, text=black!85},
  figlabel/.style={font=\sffamily\fontsize{6.4}{7.4}\selectfont,
    text=black!70, align=center},
}

% --- 表記 ---------------------------------------------------------------
% T^{b<-a} はフレーム a の座標をフレーム b へ写す同次変換。
\newcommand{\R}{\mathbb{R}}
\newcommand{\SEthree}{\mathrm{SE}(3)}
\newcommand{\SOthree}{\mathrm{SO}(3)}
\newcommand{\T}[2]{\mathbf{T}^{#1\leftarrow #2}}
\newcommand{\Rot}[2]{\mathbf{R}^{#1\leftarrow #2}}
\newcommand{\trans}[2]{\mathbf{t}^{#1\leftarrow #2}}
\newcommand{\Scene}{\mathcal{S}}
\newcommand{\Court}{\mathcal{C}}
\newcommand{\Canon}{\mathcal{N}}
\newcommand{\Cam}{\mathcal{V}}
\newcommand{\Kmat}{\mathbf{K}}
\newcommand{\point}[1]{\mathbf{x}_{#1}}
\newcommand{\pixel}[1]{\mathbf{u}_{#1}}
\newcommand{\etal}{ほか}

% --- 手法・データ・版のマクロ -------------------------------------------
\newcommand{\Method}{\textsc{GS-Court}\xspace}
\newcommand{\MethodLong}{3DGS再構成に基づくコート検出データ生成\xspace}
\newcommand{\SyntheticDataset}{\textsc{3DGS Synthetic Court V3}\xspace}
\newcommand{\BaselineName}{\textsc{TennisCourtDetector}\xspace}
\newcommand{\BaselineShort}{\textsc{TCD}\xspace}
\newcommand{\BackboneName}{\textsc{DINOv3} ViT-B/16\xspace}
\newcommand{\NHTName}{\textsc{NHT}\xspace}
\newcommand{\KPName}{\textsc{KP14}\xspace}
\newcommand{\SceneSet}{\texttt{B00}--\texttt{B03}\xspace}
\newcommand{\RunName}{\texttt{court-ms5-depth3-local-rtx-resume1}\xspace}
\newcommand{\CheckpointName}{\texttt{court-detection-epoch=17.ckpt}\xspace}
\newcommand{\BaselineCommit}{\texttt{e5cd4f1c}\xspace}
\newcommand{\ClassicBaselineCommit}{\texttt{762d0775}\xspace}
\newcommand{\TCDImageCount}{8{,}841\xspace}
\newcommand{\TCDTrainCount}{6{,}630\xspace}
\newcommand{\TCDValCount}{2{,}211\xspace}

% --- 状態ラベル ---------------------------------------------------------
\newcommand{\statusimplemented}{\textcolor{courtgreen}{\textsc{実装済}}}
\newcommand{\statuspartial}{\textcolor{poseorange}{\textsc{一部実装}}}
\newcommand{\statusunmeasured}{\textcolor{unmeasuredgray}{\textsc{未測定}}}
\newcommand{\statushypothesis}{\textcolor{poseorange}{\textsc{仮説}}}

% --- cleveref の日本語名 ------------------------------------------------
\crefname{equation}{式}{式}
\Crefname{equation}{式}{式}
\crefname{figure}{図}{図}
\Crefname{figure}{図}{図}
\crefname{table}{表}{表}
\Crefname{table}{表}{表}
\crefname{section}{節}{節}
\Crefname{section}{節}{節}
\crefname{appendix}{付録}{付録}
\Crefname{appendix}{付録}{付録}

% 実測値は results/generated/ の生成物から読み込む。
% 未測定の項目は「未測定」と明示され、値があるようには表示しない。
% =====================================================================
% 生成物: results/generated/benchmark_metrics.tex
%
% このファイルは数値の唯一の入力点である。本文・表・図はここで定義した
% マクロだけを参照し、数値を直接書かない。
%
% 出典: knowledge/runs/run-court-alignment-crossmodel-n120-s42/
%       manifest.json, metrics.json, metrics.csv, provenance.json
% manifest fingerprint:
%   bfc0285940621d6b6aa96fc79b6e308fd24e9d784341f5b7668caa964d77675c
%
% 画像を混ぜない。real_validation は公開実装付属の held-out validation
% （公式 test ではない）、synthetic_test は trajectory-group-disjoint test
% （同じ会場群の未学習軌道であり未知会場ではない）。
% =====================================================================

% 交差モデル比較は測定済みである。
% 本文・表・図は実測マクロを直接参照し、未測定の項目だけが
% \resultsUnmeasured を返す。この flag は測定状態の記録として残す。
\newif\ifresultsmeasured
\resultsmeasuredtrue

% 結果図 (figures/generated/) は実測から生成した PNG を置いてある。
\newif\ifresultsfigures
\resultsfigurestrue

% 真に未測定の項目だけがこれを返す。
\newcommand{\resultsUnmeasured}{\textcolor{unmeasuredgray}{\textbf{未測定}}}
\newcommand{\resultsPendingFigure}{\textcolor{unmeasuredgray}{\textbf{図は未生成}}}
\newcommand{\resNotApplicable}{\textcolor{unmeasuredgray}{---}}

% --- checkpoint 選択時に記録された synthetic validation -----------------
% 出典: outputs/court_detection/multiscale-depth3-local-rtx/logs/version_4
%       events.out.tfevents.*, step 29843 (epoch 17), validation は 1 回のみ。
% 注意: 同一比較条件のベンチマークではなく、実写 held-out test でもない。
\newcommand{\resCkptEpoch}{17}
\newcommand{\resCkptKPMeanPx}{1.6223}
\newcommand{\resCkptKPMedianPx}{1.2166}
\newcommand{\resCkptSegMIoU}{0.9684}
\newcommand{\resCkptLineDice}{0.7758}
\newcommand{\resCkptSemanticLineMIoU}{0.5460}
\newcommand{\resCkptPoseReprojectionPx}{21.0041}
\newcommand{\resCkptTranslationLtwoM}{0.8791}
\newcommand{\resCkptRotationGeodesicDeg}{2.6453}
\newcommand{\resCkptFocalRelativeError}{0.0609}
\newcommand{\resCkptLogFocalAbsError}{0.0620}
\newcommand{\resCkptVisiblePointCount}{7{,}346}

% --- 交差モデル比較の共通設定 -------------------------------------------
% 各 domain 決定論的 120 件、seed 42、全推論 CPU
% (CUDA_VISIBLE_DEVICES 空、torch_cuda_available=false)。
\newcommand{\resBenchSamples}{120}
\newcommand{\resBenchSeed}{42}
\newcommand{\resBenchDomainSamples}{120}
\newcommand{\resBenchLineSegmentSamples}{21}
\newcommand{\resBenchRansacThreshold}{1\%}
\newcommand{\resBenchElapsedSeconds}{410.1}
\newcommand{\resBenchManifestShort}{\texttt{bfc02859}\xspace}
\newcommand{\resBenchOursCheckpointHashShort}{\texttt{b863df1f}\xspace}

% --- real_validation / ours ---------------------------------------------
\newcommand{\resRealOursN}{120}
\newcommand{\resRealOursVisible}{1{,}673}
\newcommand{\resRealOursComplete}{1.0000}
\newcommand{\resRealOursPckHalf}{0.9797}
\newcommand{\resRealOursPckOne}{1.0000}
\newcommand{\resRealOursPckTwo}{1.0000}
\newcommand{\resRealOursPairN}{1{,}673}
\newcommand{\resRealOursPairMean}{3.1835}
\newcommand{\resRealOursPairMedian}{2.9155}
\newcommand{\resRealOursPairQ}{5.4083}
\newcommand{\resRealOursHSuccess}{1.0000}
\newcommand{\resRealOursHSuccessCount}{120/120}
\newcommand{\resRealOursLineInMean}{2.0243}
\newcommand{\resRealOursLineInMedian}{1.8563}
\newcommand{\resRealOursLineInQ}{2.8871}
\newcommand{\resRealOursLineMedian}{1.8563}
\newcommand{\resRealOursLineQ}{2.8871}
\newcommand{\resRealOursLineInFraction}{0.9988}
\newcommand{\resRealOursIouDefined}{120}
\newcommand{\resRealOursIouTotal}{120}
\newcommand{\resRealOursIouMean}{0.9878}
\newcommand{\resRealOursIouMedian}{0.9879}

% --- real_validation / tcd ----------------------------------------------
\newcommand{\resRealTcdN}{120}
\newcommand{\resRealTcdComplete}{0.9952}
\newcommand{\resRealTcdPckHalf}{0.9271}
\newcommand{\resRealTcdPckOne}{0.9779}
\newcommand{\resRealTcdPckTwo}{0.9922}
\newcommand{\resRealTcdPairN}{1{,}665}
\newcommand{\resRealTcdPairMean}{3.7136}
\newcommand{\resRealTcdPairMedian}{3.0000}
\newcommand{\resRealTcdPairQ}{6.3717}
\newcommand{\resRealTcdHSuccess}{1.0000}
\newcommand{\resRealTcdHSuccessCount}{120/120}
\newcommand{\resRealTcdLineInMean}{24.8039}
\newcommand{\resRealTcdLineInMedian}{2.1039}
\newcommand{\resRealTcdLineInQ}{3.7086}
\newcommand{\resRealTcdLineMedian}{2.1039}
\newcommand{\resRealTcdLineQ}{3.7086}
\newcommand{\resRealTcdLineInFraction}{0.9988}
\newcommand{\resRealTcdIouDefined}{119}
\newcommand{\resRealTcdIouTotal}{120}
\newcommand{\resRealTcdIouMean}{0.9858}
\newcommand{\resRealTcdIouMedian}{0.9873}
\newcommand{\resRealTcdOutlierPx}{2{,}691.3}
\newcommand{\resRealTcdOutlierVisible}{10}

% --- synthetic_test / ours ----------------------------------------------
\newcommand{\resSynthOursN}{120}
\newcommand{\resSynthOursVisible}{1{,}033}
\newcommand{\resSynthOursComplete}{1.0000}
\newcommand{\resSynthOursPckHalf}{0.7144}
\newcommand{\resSynthOursPckOne}{0.9313}
\newcommand{\resSynthOursPckTwo}{0.9816}
\newcommand{\resSynthOursPckFive}{0.9952}
\newcommand{\resSynthOursPairN}{1{,}033}
\newcommand{\resSynthOursPairMean}{7.0454}
\newcommand{\resSynthOursPairMedian}{3.6503}
\newcommand{\resSynthOursPairQ}{9.5168}
\newcommand{\resSynthOursHSuccess}{1.0000}
\newcommand{\resSynthOursHSuccessCount}{120/120}
\newcommand{\resSynthOursLineInMean}{56.7582}
\newcommand{\resSynthOursLineInMedian}{7.4264}
\newcommand{\resSynthOursLineInQ}{199.2716}
\newcommand{\resSynthOursLineInFraction}{0.7127}
\newcommand{\resSynthOursLineMedian}{422.0964}
\newcommand{\resSynthOursLineQ}{4{,}547.4852}
\newcommand{\resSynthOursIouDefined}{50}
\newcommand{\resSynthOursIouTotal}{120}
\newcommand{\resSynthOursIouMean}{0.9003}
\newcommand{\resSynthOursIouMedian}{0.9490}

% synthetic_test / ours の coverage 層別（主指標を層で割った内訳）
\newcommand{\resSynthOursFullN}{7}
\newcommand{\resSynthOursFullLineMedian}{4.2441}
\newcommand{\resSynthOursFullPckOne}{0.9490}
\newcommand{\resSynthOursFullVisible}{98}
\newcommand{\resSynthOursNearFullN}{53}
\newcommand{\resSynthOursNearFullLineMedian}{6.3320}
\newcommand{\resSynthOursNearFullPckOne}{0.9318}
\newcommand{\resSynthOursNearFullVisible}{557}
\newcommand{\resSynthOursPartialN}{60}
\newcommand{\resSynthOursPartialLineMedian}{11.3766}
\newcommand{\resSynthOursPartialPckOne}{0.9259}
\newcommand{\resSynthOursPartialVisible}{378}

% synthetic_test / ours の会場別 PCK@0.01
\newcommand{\resSynthOursSceneBzeroN}{28}
\newcommand{\resSynthOursSceneBzeroPckOne}{0.9600}
\newcommand{\resSynthOursSceneBzeroVisible}{250}
\newcommand{\resSynthOursSceneBoneN}{18}
\newcommand{\resSynthOursSceneBonePckOne}{0.7452}
\newcommand{\resSynthOursSceneBoneVisible}{157}
\newcommand{\resSynthOursSceneBtwoN}{38}
\newcommand{\resSynthOursSceneBtwoPckOne}{0.9829}
\newcommand{\resSynthOursSceneBtwoVisible}{351}
\newcommand{\resSynthOursSceneBthreeN}{36}
\newcommand{\resSynthOursSceneBthreePckOne}{0.9455}
\newcommand{\resSynthOursSceneBthreeVisible}{275}

% --- synthetic_test / tcd -----------------------------------------------
\newcommand{\resSynthTcdN}{120}
\newcommand{\resSynthTcdComplete}{0.0281}
\newcommand{\resSynthTcdPckHalf}{0.0010}
\newcommand{\resSynthTcdPckOne}{0.0019}
\newcommand{\resSynthTcdPckTwo}{0.0058}
\newcommand{\resSynthTcdPckFive}{0.0126}
\newcommand{\resSynthTcdPairN}{29}
\newcommand{\resSynthTcdPairMean}{108.5576}
\newcommand{\resSynthTcdPairMedian}{66.8262}
\newcommand{\resSynthTcdPairQ}{178.8752}
\newcommand{\resSynthTcdHSuccess}{0.0167}
\newcommand{\resSynthTcdHSuccessCount}{2/120}
\newcommand{\resSynthTcdLineInMean}{156.2766}
\newcommand{\resSynthTcdLineInMedian}{156.2766}
\newcommand{\resSynthTcdLineInQ}{172.5752}
\newcommand{\resSynthTcdLineInFraction}{0.0124}
\newcommand{\resSynthTcdLineMedian}{1{,}517.8520}
\newcommand{\resSynthTcdLineQ}{2{,}569.1827}
\newcommand{\resSynthTcdIouDefined}{0}
\newcommand{\resSynthTcdIouTotal}{120}
\newcommand{\resSynthTcdInsufficientFailures}{117}
\newcommand{\resSynthTcdRansacFailures}{1}
\newcommand{\resSynthTcdSceneBzeroN}{28}
\newcommand{\resSynthTcdSceneBzeroComplete}{0.1160}
\newcommand{\resSynthTcdSceneBzeroHSuccess}{2}
\newcommand{\resSynthTcdSceneBzeroPairs}{29}
\newcommand{\resSynthTcdSceneOtherValid}{0}

% --- 古典 baseline 補助 stress probe ------------------------------------
% 主比較とは別の決定論的な 10 real + 10 B00 synthetic 標本である。
% 主表の数値と直接比較しない。
\newcommand{\resClassicCommitShort}{\texttt{762d0775}\xspace}
\newcommand{\resClassicPatchHashShort}{\texttt{12422aec}\xspace}
\newcommand{\resClassicEnv}{Ubuntu 22.04 / OpenCV 4.5.4 / CPU のみの Docker\xspace}
\newcommand{\resClassicProbeSamples}{real 10 件 + \texttt{B00} synthetic 10 件}
\newcommand{\resClassicRealSuccess}{10/10}
\newcommand{\resClassicRealFinitePairs}{16}
\newcommand{\resClassicRealInFramePartial}{1/10}
\newcommand{\resClassicRealInFramePairs}{8/16}
\newcommand{\resClassicSynthSuccess}{6/10}
\newcommand{\resClassicSynthFailure}{4/10}
\newcommand{\resClassicSynthCrash}{3}
\newcommand{\resClassicSynthExitThree}{1}
\newcommand{\resClassicBzeroTestFrames}{238}
\newcommand{\resClassicBzeroCourts}{2}

% --- データセット集計 (manifest から確認済み) ---------------------------
\newcommand{\datasetAcceptedTotal}{8{,}415}
\newcommand{\datasetAcceptedBzero}{2{,}176}
\newcommand{\datasetAcceptedBone}{2{,}052}
\newcommand{\datasetAcceptedBtwo}{2{,}094}
\newcommand{\datasetAcceptedBthree}{2{,}093}
\newcommand{\datasetTrainFrames}{6{,}594}
\newcommand{\datasetValidationFrames}{935}
\newcommand{\datasetTestFrames}{886}
\newcommand{\datasetTrainGroups}{113}
\newcommand{\datasetValidationGroups}{15}
\newcommand{\datasetTestGroups}{13}
\newcommand{\datasetEarlierReportTotal}{8{,}433}

